\documentclass[12pt]{article}
\usepackage[margin=1in]{geometry}
\usepackage{amsmath,amssymb,amsthm}
\usepackage{hyperref}
\usepackage{enumitem}
\usepackage{booktabs}
\usepackage{caption}

\theoremstyle{plain}
\newtheorem{theorem}{Theorem}[section]
\newtheorem{proposition}[theorem]{Proposition}
\newtheorem{lemma}[theorem]{Lemma}
\newtheorem{corollary}[theorem]{Corollary}
\theoremstyle{definition}
\newtheorem{definition}[theorem]{Definition}
\newtheorem{remark}[theorem]{Remark}
\numberwithin{equation}{section}

\newcommand{\R}{\mathbb{R}}
\newcommand{\C}{\mathbb{C}}
\newcommand{\Z}{\mathbb{Z}}
\newcommand{\HH}{\mathbb{H}}
\newcommand{\br}{\mathrm{br}}
\newcommand{\vis}{\mathrm{vis}}
\newcommand{\eff}{\mathrm{eff}}
\newcommand{\LoN}{\mathrm{LoN}}
\newcommand{\act}{\mathrm{act}}
\newcommand{\ext}{\mathrm{ext}}
\newcommand{\stat}{\mathrm{stat}}
\newcommand{\well}{\mathrm{well}}
\newcommand{\bari}{\mathrm{bar}}
\newcommand{\fU}{\mathfrak{U}}
\newcommand{\fI}{\mathfrak{I}}
\newcommand{\fV}{\mathfrak{V}}
\newcommand{\fB}{\mathfrak{B}}
\DeclareMathOperator{\Tr}{Tr}
\DeclareMathOperator{\Ran}{Ran}
\DeclareMathOperator{\Cov}{Cov}
\DeclareMathOperator{\sech}{sech}
\DeclareMathOperator{\rank}{rank}

\begin{document}

\begin{titlepage}
\centering
\vspace*{0.8cm}

{\LARGE\bfseries Quantum Chemistry from LoNalogy:\\[6pt]
Canonical Bridge Selection,\\[6pt]
Exact Feshbach--Schur Downfolding,\\[6pt]
and Bridge Covariance Functionals}

\vspace{1.2cm}

{\large Masayoshi Nakamura$^{1,*}$}

\vspace{0.6cm}

{\normalsize
$^1$\,Kitayono Mental Clinic,
338-0002, 4-20-14 Shimoochiai, Chuo-ku,\\
Saitama City, Saitama, Japan.
\texttt{mjolnir501@gmail.com}\\[4pt]
$^*$\,Corresponding author.
}

\vspace{0.6cm}

{\small
\noindent\textbf{Related deposits:}\\
LoNalogy framework (v87.1):
\href{https://doi.org/10.5281/zenodo.19115338}{doi:10.5281/zenodo.19115338}\\
Yang--Mills proof (v90):
\href{https://doi.org/10.5281/zenodo.19183838}{doi:10.5281/zenodo.19183838}\\
Companion papers: Statistical mechanics on $X(2)$,
Particle cosmology (v9), Hubble reconstruction (v2),
Coronal heating, Extra dimensions
}

\vspace{0.8cm}

\begin{abstract}
\noindent
We apply the LoNalogy three-sector framework to quantum chemistry.
Three results are established with complete proofs:

\textbf{(A) Canonical bridge selection and exact downfolding.}
Given any active-space projector $P_A$, the bridge projector
$P_B:=\Pi_{\Ran(QHP_A)}$ is canonically defined, yielding an exact
three-sector Hamiltonian with $P_A H P_I=0$.  The exact
energy-dependent effective Hamiltonian
$H_A^{\eff}(E)=H_{AA}-H_{AB}(H_{BB}-E-\Sigma_B(E))^{-1}H_{BA}$
is the algebraic parent of CASPT2, NEVPT2, and DUCC.
A static generalised eigenproblem closure is derived.

\textbf{(B) System-dependent bridge covariance functional.}
The exact correlation correction from hidden-mode elimination is
the path-integrated bridge covariance
$E_{\br,\mathrm{cov}}[X|X_0]
=-\int_0^1(1-s)\langle\delta X,
\mathcal{C}_{\br}[X_s]\delta X\rangle\,ds$,
where
$\mathcal{C}_{\br}[X]=\beta\,\Cov_X(\Phi_X,\Phi_X)\succeq 0$.
A universal fixed-kernel quadratic $E_{xc}$ is shown to be
incompatible with the derivative discontinuity of exact DFT (no-go).
In the Gaussianised bridge regime, the kernel reduces exactly to
$MD^{-1}M^\top$.

\textbf{(C) Modified P\"oschl--Teller local normal form.}
At any non-degenerate stationary point of a 1D reaction path,
a canonical coordinate change brings the PES into $\sech^2$ form
with agreement through fourth order.  A no-go theorem shows this
cannot be a universal PES law: the obstruction is the cubic
anharmonicity $V'''(q_*)\neq 0$ of generic potentials.

Four no-go theorems delineate the exact boundaries of applicability.
\end{abstract}
\end{titlepage}

\setcounter{page}{1}
\tableofcontents
\newpage


%% ============================================================
\section*{Introduction}
\addcontentsline{toc}{section}{Introduction}
%% ============================================================

Quantum chemistry faces a hierarchy of computational challenges:
the mean-field approximation (Hartree--Fock) misses instantaneous
electron correlation; density functional theory (DFT) is in
principle exact but depends on an unknown exchange-correlation
functional $E_{xc}[\rho]$; post-Hartree--Fock methods (CI, CC,
DMRG) are systematically improvable but scale poorly; and
fermionic quantum Monte Carlo suffers from the sign problem.

The LoNalogy framework~\cite{LoNalogyv87,YMmassGap} was not
designed for quantum chemistry.  Its core---the three-sector
decomposition $\fU=\fI\oplus\fV\oplus\fB$, the Schur complement
(Kimura--Th\'evenin principle), the $\sech^2$ bridge weight, and
the projected determinant bundle---was developed for quantum field
theory and the Yang--Mills mass gap.  However, the mathematical
structures are universal enough to apply to the multi-electron
problem.

This paper identifies exactly where LoNalogy tools close quantum
chemistry problems and where they provably do not (no-go).
The results are organised in three parts, each containing both
positive theorems and no-go results.


%% ============================================================
\part{Canonical Bridge Selection and Exact Downfolding}
%% ============================================================

\section{Canonical Three-Sector Decomposition}\label{sec:canonical}

\begin{definition}[Active-space projector]\label{def:active}
Let $\mathcal{H}$ be a finite-dimensional many-electron Hilbert
space and $H:\mathcal{H}\to\mathcal{H}$ a self-adjoint Hamiltonian.
Let $P_A:\mathcal{H}\to\mathcal{H}_A$ be an orthogonal projector
onto an \emph{active subspace} $\mathcal{H}_A$ of dimension $n_A$.
Define $Q:=I-P_A$ (external projector) and
$\mathcal{H}_Q:=Q\mathcal{H}$.
\end{definition}

\begin{definition}[Canonical bridge projector]\label{def:bridge}
The \emph{canonical bridge projector} is
\begin{equation}\label{eq:PB}
\boxed{P_B := \Pi_{\Ran(QHP_A)},}
\end{equation}
i.e., the orthogonal projector onto the range of the operator
$QHP_A:\mathcal{H}_A\to\mathcal{H}_Q$.  The \emph{idea (core)
projector} is
\begin{equation}
P_I := I - P_A - P_B = Q - P_B.
\end{equation}
The three subspaces satisfy
$\mathcal{H}=\mathcal{H}_A\oplus\mathcal{H}_B\oplus\mathcal{H}_I$
with $\mathcal{H}_B:=P_B\mathcal{H}$ and
$\mathcal{H}_I:=P_I\mathcal{H}$.
\end{definition}

\begin{remark}[Why canonical?]
The bridge is not chosen by hand.  Given any active space, the
bridge is uniquely determined as the image of $QHP_A$: it consists
of exactly those external states that are \emph{directly coupled}
to the active space by $H$.  States orthogonal to this image have
no direct matrix element with the active space.
\end{remark}

\begin{theorem}[Canonical decoupling]\label{thm:decoupling}
With $P_B$ defined by~\eqref{eq:PB}:
\begin{equation}\label{eq:decoupling}
\boxed{P_A H P_I = 0 = P_I H P_A.}
\end{equation}
Consequently, the Hamiltonian has the three-sector block form
\begin{equation}\label{eq:blockH}
H = \begin{pmatrix}
H_{AA} & H_{AB} & 0\\
H_{BA} & H_{BB} & H_{BI}\\
0 & H_{IB} & H_{II}
\end{pmatrix}
\end{equation}
with respect to the decomposition
$\mathcal{H}=\mathcal{H}_A\oplus\mathcal{H}_B\oplus\mathcal{H}_I$.
\end{theorem}

\begin{proof}
\emph{Step 1: Show $P_I H P_A = 0$.}

By definition, $P_B$ is the orthogonal projector onto
$\Ran(QHP_A)$.  Therefore $P_I = Q - P_B$ projects onto the
orthogonal complement of $\Ran(QHP_A)$ within $\mathcal{H}_Q$.

For any $|\psi_A\rangle\in\mathcal{H}_A$, the vector
$QH|\psi_A\rangle = QHP_A|\psi_A\rangle$ lies in $\Ran(QHP_A)$
by definition.  Therefore
\begin{equation}
P_I(QHP_A|\psi_A\rangle) = 0
\end{equation}
since $P_I$ projects onto $\Ran(QHP_A)^\perp\cap\mathcal{H}_Q$.

Now, $P_IHP_A = P_IQHP_A$ because $P_IP_A=0$ (the projectors
are mutually orthogonal) and $P_IQ=P_I$ (since $P_I\leq Q$).
Therefore $P_IHP_A = P_I(QHP_A) = 0$.

\emph{Step 2: Show $P_AHP_I = 0$.}

Since $H$ is self-adjoint:
\begin{equation}
P_AHP_I = (P_IHP_A)^\dagger = 0^\dagger = 0.
\end{equation}

\emph{Step 3: Block form.}

The off-diagonal blocks of $H$ in the three-sector decomposition
are:
\begin{align}
H_{AI} &= P_AHP_I = 0,\\
H_{IA} &= P_IHP_A = 0,\\
H_{AB} &= P_AHP_B \quad(\text{generally }\neq 0),\\
H_{BI} &= P_BHP_I \quad(\text{generally }\neq 0).
\end{align}
This gives~\eqref{eq:blockH}.
\end{proof}

\begin{proposition}[Bridge rank bound]\label{prop:rankbound}
The bridge dimension satisfies
\begin{equation}
\dim\mathcal{H}_B = \rank(QHP_A) \leq
\min(\dim\mathcal{H}_A,\,\dim\mathcal{H}_Q).
\end{equation}
In particular, $\dim\mathcal{H}_B\leq\dim\mathcal{H}_A=n_A$.
\end{proposition}

\begin{proof}
The rank of $QHP_A:\mathcal{H}_A\to\mathcal{H}_Q$ is at most
$\min(\dim\mathcal{H}_A,\dim\mathcal{H}_Q)$.  Since
$\dim\mathcal{H}_B=\dim\Ran(QHP_A)=\rank(QHP_A)$, the bound
follows.
\end{proof}

\begin{remark}[Genuine reduction]
The bridge rank is at most $n_A$, the active-space dimension.
In contrast, the full external space $\mathcal{H}_Q$ can be
exponentially large.  The canonical bridge captures \emph{only}
those external states directly coupled to the active space,
discarding the vast $\mathcal{H}_I$ component that communicates
with the active space only through the bridge.
\end{remark}


\section{Exact Feshbach--Schur Downfolding}\label{sec:feshbach}

\begin{theorem}[Exact three-sector downfolding]\label{thm:downfolding}
Define $K(E):=H-EI$.  Suppose $H_{II}-E$ is invertible and
\begin{equation}\label{eq:Htilde}
\widetilde{H}_B(E) := H_{BB}-E-H_{BI}(H_{II}-E)^{-1}H_{IB}
\end{equation}
is also invertible.  Then the eigenvalue equation $H\Psi=E\Psi$
is equivalent to the active-space equation
\begin{equation}\label{eq:activeEV}
H_A^{\eff}(E)\,\psi_A = E\,\psi_A,
\end{equation}
where $\psi_A=P_A\Psi$ and the exact effective Hamiltonian is
\begin{equation}\label{eq:Heff}
\boxed{H_A^{\eff}(E) = H_{AA}
- H_{AB}\,\widetilde{H}_B(E)^{-1}\,H_{BA}.}
\end{equation}
Equivalently, defining the bridge self-energy
$\Sigma_B(E):=H_{BI}(H_{II}-E)^{-1}H_{IB}$:
\begin{equation}\label{eq:Heff2}
\boxed{H_A^{\eff}(E) = H_{AA}
- H_{AB}\bigl(H_{BB}-E-\Sigma_B(E)\bigr)^{-1}H_{BA}.}
\end{equation}
\end{theorem}

\begin{proof}
\emph{Step 1: Write the block eigenvalue equation.}

Using the block form~\eqref{eq:blockH}, the equation
$H\Psi=E\Psi$ with $\Psi=(\psi_A,\psi_B,\psi_I)^\top$ gives:
\begin{align}
H_{AA}\psi_A + H_{AB}\psi_B &= E\psi_A,\label{eq:block1}\\
H_{BA}\psi_A + H_{BB}\psi_B + H_{BI}\psi_I &= E\psi_B,\label{eq:block2}\\
H_{IB}\psi_B + H_{II}\psi_I &= E\psi_I.\label{eq:block3}
\end{align}

\emph{Step 2: Eliminate $\psi_I$.}

From~\eqref{eq:block3}:
\begin{equation}
(H_{II}-E)\psi_I = -H_{IB}\psi_B.
\end{equation}
Since $H_{II}-E$ is invertible by hypothesis:
\begin{equation}\label{eq:psiI}
\psi_I = -(H_{II}-E)^{-1}H_{IB}\psi_B.
\end{equation}

\emph{Step 3: Substitute into~\eqref{eq:block2}.}

Inserting~\eqref{eq:psiI} into~\eqref{eq:block2}:
\begin{align}
H_{BA}\psi_A + H_{BB}\psi_B
- H_{BI}(H_{II}-E)^{-1}H_{IB}\psi_B &= E\psi_B,\\
H_{BA}\psi_A
+ \bigl(H_{BB}-E-\underbrace{H_{BI}(H_{II}-E)^{-1}H_{IB}}_{\Sigma_B(E)}\bigr)\psi_B &= 0.
\end{align}
Therefore
\begin{equation}
\widetilde{H}_B(E)\psi_B = -H_{BA}\psi_A.
\end{equation}
Since $\widetilde{H}_B(E)$ is invertible:
\begin{equation}\label{eq:psiB}
\psi_B = -\widetilde{H}_B(E)^{-1}H_{BA}\psi_A.
\end{equation}

\emph{Step 4: Substitute into~\eqref{eq:block1}.}

Inserting~\eqref{eq:psiB} into~\eqref{eq:block1}:
\begin{align}
H_{AA}\psi_A - H_{AB}\widetilde{H}_B(E)^{-1}H_{BA}\psi_A
&= E\psi_A,\\
\bigl(H_{AA}-H_{AB}\widetilde{H}_B(E)^{-1}H_{BA}\bigr)\psi_A
&= E\psi_A.
\end{align}
This is~\eqref{eq:activeEV} with $H_A^{\eff}(E)$ as
in~\eqref{eq:Heff}.

\emph{Step 5: Verify equivalence.}

Given $\psi_A$ satisfying~\eqref{eq:activeEV}, we can reconstruct
$\psi_B$ via~\eqref{eq:psiB} and $\psi_I$ via~\eqref{eq:psiI}.
Then $\Psi=(\psi_A,\psi_B,\psi_I)^\top$ satisfies $H\Psi=E\Psi$.
Conversely, any eigenvector $\Psi$ of $H$ with eigenvalue $E$
(outside the spectrum of $H_{II}$ and $\widetilde{H}_B(E)$)
yields $\psi_A$ satisfying~\eqref{eq:activeEV}.
\end{proof}

\begin{remark}[Relation to standard methods]
The effective Hamiltonian~\eqref{eq:Heff2} is the algebraic parent
of several standard quantum chemistry methods:
CASPT2 approximates $\Sigma_B(E)$ by second-order perturbation
theory;
NEVPT2 uses a different zeroth-order Hamiltonian;
DUCC (double unitary coupled cluster) constructs an
energy-independent downfolded Hamiltonian by absorbing external
amplitudes.
All are approximations to the exact $H_A^{\eff}(E)$.
\end{remark}


\section{Static Closure}\label{sec:static}

\begin{theorem}[Fixed-point iteration]\label{thm:fixedpoint}
Define $F:E\mapsto\lambda_0(H_A^{\eff}(E))$, where $\lambda_0$
denotes the lowest eigenvalue.  If there exists a closed interval
$I\subset\R$ such that $F:I\to I$ and
\begin{equation}
\sup_{E\in I}\left|\frac{\partial F}{\partial E}(E)\right| < 1,
\end{equation}
then the iteration $E^{(n+1)}:=F(E^{(n)})$ converges to a unique
fixed point $E_*\in I$ satisfying
$E_*=\lambda_0(H_A^{\eff}(E_*))$.
\end{theorem}

\begin{proof}
This is the Banach contraction mapping theorem applied to $F$ on
the complete metric space $(I,|\cdot|)$.  The contraction condition
$|F(E_1)-F(E_2)|\leq c|E_1-E_2|$ with $c<1$ follows from the
mean value theorem and the bound on $|\partial F/\partial E|$.
\end{proof}

\begin{theorem}[Static generalised eigenproblem]\label{thm:static}
Fix a reference energy $E_*$ and define
\begin{align}
R_* &:= \left.\frac{\partial H_A^{\eff}}{\partial E}\right|_{E=E_*},\\
H_*^{\stat} &:= H_A^{\eff}(E_*) - E_*\,R_*,\\
S_* &:= I - R_*.
\end{align}
Then the exact fixed-point equation
$H_A^{\eff}(E)\psi_A=E\psi_A$ is approximated to first order in
$(E-E_*)$ by the generalised eigenproblem
\begin{equation}\label{eq:geneig}
\boxed{H_*^{\stat}\,\psi_A = E\,S_*\,\psi_A.}
\end{equation}
\end{theorem}

\begin{proof}
Expand $H_A^{\eff}(E)$ to first order around $E_*$:
\begin{equation}
H_A^{\eff}(E) \approx H_A^{\eff}(E_*) + (E-E_*)R_*.
\end{equation}
The eigenvalue equation $H_A^{\eff}(E)\psi_A=E\psi_A$ becomes
\begin{align}
\bigl(H_A^{\eff}(E_*)+(E-E_*)R_*\bigr)\psi_A &= E\psi_A,\\
H_A^{\eff}(E_*)\psi_A + ER_*\psi_A - E_*R_*\psi_A &= E\psi_A,\\
\bigl(H_A^{\eff}(E_*)-E_*R_*\bigr)\psi_A &= E(I-R_*)\psi_A.
\end{align}
This is~\eqref{eq:geneig}.
\end{proof}

\begin{proposition}[Energy derivative of $H_A^{\eff}$]\label{prop:Rstar}
The energy derivative of the effective Hamiltonian is
\begin{equation}\label{eq:Rstar}
R_* = \frac{\partial H_A^{\eff}}{\partial E}\bigg|_{E_*}
= H_{AB}\,\widetilde{H}_B(E_*)^{-1}\,
\frac{\partial\widetilde{H}_B}{\partial E}\bigg|_{E_*}\,
\widetilde{H}_B(E_*)^{-1}\,H_{BA},
\end{equation}
where
\begin{equation}
\frac{\partial\widetilde{H}_B}{\partial E}
= -I_B - H_{BI}(H_{II}-E)^{-2}H_{IB}.
\end{equation}
\end{proposition}

\begin{proof}
Differentiating~\eqref{eq:Heff} with respect to $E$:
\begin{equation}
\frac{\partial H_A^{\eff}}{\partial E}
= -H_{AB}\,\frac{\partial}{\partial E}
\bigl(\widetilde{H}_B(E)^{-1}\bigr)\,H_{BA}.
\end{equation}
Using the identity
$\frac{d}{dE}(M(E)^{-1})=-M^{-1}\frac{dM}{dE}M^{-1}$:
\begin{equation}
\frac{\partial H_A^{\eff}}{\partial E}
= H_{AB}\,\widetilde{H}_B^{-1}\,
\frac{\partial\widetilde{H}_B}{\partial E}\,
\widetilde{H}_B^{-1}\,H_{BA}.
\end{equation}
From~\eqref{eq:Htilde}:
$\frac{\partial\widetilde{H}_B}{\partial E}
=\frac{\partial}{\partial E}(H_{BB}-E-\Sigma_B(E))
=-I_B-\frac{\partial\Sigma_B}{\partial E}$,
and $\frac{\partial\Sigma_B}{\partial E}
=H_{BI}(H_{II}-E)^{-2}H_{IB}$.
\end{proof}


%% ============================================================
\part{Bridge Covariance Functional}
%% ============================================================

\section{Log-Marginal Hessian}\label{sec:hessian}

\begin{definition}[Reduced free energy]\label{def:Gamma}
Let $X\in\R^d$ be an active descriptor (density, 1-RDM, or any
set of active-space parameters) and $Z\in\R^m$ be hidden/external
variables.  Let $\Phi(X,Z)$ be the full energy functional.
The reduced free energy at inverse temperature $\beta>0$ is
\begin{equation}\label{eq:Gamma}
\Gamma_\beta(X) := -\beta^{-1}\log\int_{\R^m}
e^{-\beta\Phi(X,Z)}\,dZ.
\end{equation}
\end{definition}

\begin{theorem}[Exact bridge covariance Hessian]\label{thm:hessian}
If $\Gamma_\beta$ is $C^2$ at $X$, then
\begin{equation}\label{eq:hessian}
\boxed{\nabla_X^2\Gamma_\beta(X)
= \mathbb{E}_X[\Phi_{XX}]
- \beta\,\Cov_X(\Phi_X,\Phi_X),}
\end{equation}
where $\Phi_{XX}:=\nabla_X^2\Phi(X,Z)$,
$\Phi_X:=\nabla_X\Phi(X,Z)$, and the expectation and covariance
are taken with respect to the conditional Gibbs measure
\begin{equation}
\mu_X(dZ) = \frac{e^{-\beta\Phi(X,Z)}}
{\int e^{-\beta\Phi(X,Z')}\,dZ'}\,dZ.
\end{equation}
\end{theorem}

\begin{proof}
\emph{Step 1: First derivative.}

Differentiating~\eqref{eq:Gamma}:
\begin{align}
\nabla_X\Gamma_\beta(X)
&= -\beta^{-1}\nabla_X\log\int e^{-\beta\Phi}\,dZ\\
&= -\beta^{-1}\frac{\int(-\beta\Phi_X)e^{-\beta\Phi}\,dZ}
{\int e^{-\beta\Phi}\,dZ}\\
&= \frac{\int\Phi_X\,e^{-\beta\Phi}\,dZ}
{\int e^{-\beta\Phi}\,dZ}\\
&= \mathbb{E}_X[\Phi_X].
\end{align}

\emph{Step 2: Second derivative.}

Differentiating $\nabla_X\Gamma_\beta=\mathbb{E}_X[\Phi_X]$
with respect to $X$.  Write
$\langle f\rangle:=\mathbb{E}_X[f]$ for brevity.
\begin{equation}
\nabla_X\langle\Phi_X\rangle
= \nabla_X\!\left(\frac{\int\Phi_X\,e^{-\beta\Phi}\,dZ}
{\int e^{-\beta\Phi}\,dZ}\right).
\end{equation}
Using the quotient rule for $\nabla_X(N/D)=(D\nabla_X N-N\nabla_X D)/D^2$
with $N=\int\Phi_X e^{-\beta\Phi}\,dZ$ and $D=\int e^{-\beta\Phi}\,dZ$:
\begin{align}
\nabla_X N &= \int\bigl(\Phi_{XX}-\beta\Phi_X\Phi_X^\top\bigr)
e^{-\beta\Phi}\,dZ,\\
\nabla_X D &= \int(-\beta\Phi_X)e^{-\beta\Phi}\,dZ
= -\beta\,D\,\langle\Phi_X\rangle.
\end{align}
Therefore:
\begin{align}
\nabla_X^2\Gamma_\beta
&= \frac{D\bigl(\int(\Phi_{XX}-\beta\Phi_X\Phi_X^\top)
e^{-\beta\Phi}dZ\bigr)
-N\bigl(-\beta D\langle\Phi_X\rangle\bigr)}{D^2}\\
&= \langle\Phi_{XX}\rangle
-\beta\langle\Phi_X\Phi_X^\top\rangle
+\beta\langle\Phi_X\rangle\langle\Phi_X\rangle^\top\\
&= \langle\Phi_{XX}\rangle
-\beta\bigl(\langle\Phi_X\Phi_X^\top\rangle
-\langle\Phi_X\rangle\langle\Phi_X\rangle^\top\bigr)\\
&= \mathbb{E}_X[\Phi_{XX}]-\beta\,\Cov_X(\Phi_X,\Phi_X).
\end{align}
This is~\eqref{eq:hessian}.
\end{proof}


\section{Bridge Covariance Kernel and Functional}\label{sec:covariance}

\begin{definition}[Bridge covariance kernel]\label{def:Cbr}
The \emph{system-dependent bridge covariance kernel} is
\begin{equation}\label{eq:Cbr}
\boxed{\mathcal{C}_{\br}[X]
:= \beta\,\Cov_X(\Phi_X,\Phi_X) \succeq 0.}
\end{equation}
This is a positive semidefinite $d\times d$ matrix for each $X$.
\end{definition}

\begin{theorem}[Exact path-integrated bridge covariance functional]\label{thm:functional}
Fix a reference point $X_0$.  Let $\delta X:=X-X_0$ and
$X_s:=X_0+s\,\delta X$ for $s\in[0,1]$.  If $\Gamma_\beta$ is
$C^2$ on the line segment $\{X_s:s\in[0,1]\}$, then
\begin{equation}
\Gamma_\beta(X) = \Gamma_\beta(X_0)
+ \langle\nabla\Gamma_\beta(X_0),\delta X\rangle
+ E_{\mathrm{bare}}^{(2)}[X|X_0]
+ E_{\br,\mathrm{cov}}[X|X_0],
\end{equation}
where
\begin{align}
E_{\mathrm{bare}}^{(2)}[X|X_0]
&:= \int_0^1(1-s)\,
\langle\delta X,\,\mathbb{E}_{X_s}[\Phi_{XX}]\,\delta X\rangle\,ds,\\
\boxed{E_{\br,\mathrm{cov}}[X|X_0]
:= -\int_0^1(1-s)\,
\langle\delta X,\,\mathcal{C}_{\br}[X_s]\,\delta X\rangle\,ds.}
\end{align}
\end{theorem}

\begin{proof}
By Taylor's theorem with integral remainder:
\begin{equation}
\Gamma_\beta(X) = \Gamma_\beta(X_0)
+\langle\nabla\Gamma_\beta(X_0),\delta X\rangle
+\int_0^1(1-s)\langle\delta X,
\nabla_X^2\Gamma_\beta(X_s)\delta X\rangle\,ds.
\end{equation}
Substituting the Hessian decomposition~\eqref{eq:hessian}:
\begin{align}
\int_0^1&(1-s)\langle\delta X,
\nabla_X^2\Gamma_\beta(X_s)\delta X\rangle\,ds\\
&= \int_0^1(1-s)\langle\delta X,
\mathbb{E}_{X_s}[\Phi_{XX}]\delta X\rangle\,ds
-\int_0^1(1-s)\langle\delta X,
\mathcal{C}_{\br}[X_s]\delta X\rangle\,ds\\
&= E_{\mathrm{bare}}^{(2)}[X|X_0]
+ E_{\br,\mathrm{cov}}[X|X_0].
\end{align}
\end{proof}


\section{Gaussianised Bridge Regime}\label{sec:gaussian}

\begin{theorem}[Exact quadratic kernel for Gaussian bridge]\label{thm:gaussian}
Suppose the full energy has the form
\begin{equation}\label{eq:gaussenergy}
\Phi(X,Z) = \Phi_0(X) + \tfrac{1}{2}Z^\top DZ - X^\top MZ,
\end{equation}
where $D\succ 0$ is a fixed positive definite matrix and
$M\in\R^{d\times m}$.  Then the Gaussian integral over $Z$ gives
\begin{equation}\label{eq:gaussGamma}
\Gamma_\beta(X) = \Phi_0(X)
-\tfrac{1}{2}X^\top MD^{-1}M^\top X
+\frac{1}{2\beta}\log\det D + \mathrm{const},
\end{equation}
and the bridge covariance kernel is
\begin{equation}\label{eq:gausskernel}
\boxed{\mathcal{C}_{\br}[X] \equiv MD^{-1}M^\top}
\end{equation}
independent of $X$.  The bridge covariance functional becomes
exactly quadratic:
\begin{equation}\label{eq:gaussfunctional}
E_{\br,\mathrm{cov}}[X|X_0]
= -\tfrac{1}{2}\langle X-X_0,\,MD^{-1}M^\top\,(X-X_0)\rangle.
\end{equation}
\end{theorem}

\begin{proof}
\emph{Step 1: Complete the square.}

From~\eqref{eq:gaussenergy}:
\begin{align}
\Phi(X,Z) &= \Phi_0(X)+\tfrac{1}{2}Z^\top DZ-X^\top MZ\\
&= \Phi_0(X)+\tfrac{1}{2}(Z-D^{-1}M^\top X)^\top
D(Z-D^{-1}M^\top X)
-\tfrac{1}{2}X^\top MD^{-1}M^\top X.
\end{align}

\emph{Step 2: Gaussian integral.}

\begin{align}
\int e^{-\beta\Phi(X,Z)}\,dZ
&= e^{-\beta\Phi_0(X)+\frac{\beta}{2}X^\top MD^{-1}M^\top X}
\int e^{-\frac{\beta}{2}(Z-\cdots)^\top D(Z-\cdots)}\,dZ\\
&= e^{-\beta\Phi_0(X)+\frac{\beta}{2}X^\top MD^{-1}M^\top X}
\cdot\left(\frac{(2\pi)^m}{\beta^m\det D}\right)^{1/2}.
\end{align}

\emph{Step 3: Take log.}

\begin{align}
\Gamma_\beta(X) &= -\beta^{-1}\log(\cdots)\\
&= \Phi_0(X)-\tfrac{1}{2}X^\top MD^{-1}M^\top X
+\frac{1}{2\beta}\log\det D+\text{const}.
\end{align}

\emph{Step 4: Compute Hessian.}

$\nabla_X^2\Gamma_\beta=\nabla_X^2\Phi_0-MD^{-1}M^\top$.
But $\mathbb{E}_X[\Phi_{XX}]=\nabla_X^2\Phi_0$
(since $\Phi_{XX}=(\Phi_0)_{XX}$ is independent of $Z$).
Therefore $\mathcal{C}_{\br}[X]=\mathbb{E}_X[\Phi_{XX}]-\nabla_X^2\Gamma_\beta=MD^{-1}M^\top$.

\emph{Step 5: Functional.}

$E_{\br,\mathrm{cov}}[X|X_0]
=-\int_0^1(1-s)\langle\delta X,MD^{-1}M^\top\delta X\rangle\,ds
=-\frac{1}{2}\langle\delta X,MD^{-1}M^\top\delta X\rangle$.
\end{proof}


\section{No-Go: Universal Fixed-Kernel $E_{xc}$}\label{sec:nogo_Exc}

\begin{theorem}[No-go for universal fixed quadratic $E_{xc}$]\label{thm:nogo_Exc}
There is no fixed self-adjoint operator $\Delta$ and constant
$\eta>0$ such that
\begin{equation}
E_{xc}[\rho] = \eta\langle\rho,\Delta\rho\rangle
\end{equation}
is the exact universal exchange-correlation functional for all
electron numbers $N$.
\end{theorem}

\begin{proof}
\emph{Step 1: Smoothness of the quadratic form.}

If $E_{xc}[\rho]=\eta\langle\rho,\Delta\rho\rangle$ with fixed
$\Delta$, then
\begin{equation}
v_{xc}[\rho]=\frac{\delta E_{xc}}{\delta\rho}=2\eta\Delta\rho,
\qquad
f_{xc}[\rho]=\frac{\delta^2 E_{xc}}{\delta\rho\,\delta\rho}
=2\eta\Delta.
\end{equation}
Both are everywhere smooth (infinitely differentiable) in $\rho$
for any fixed bounded $\Delta$.

\emph{Step 2: Derivative discontinuity requirement.}

The exact $E_{xc}$ of Kohn--Sham DFT exhibits a derivative
discontinuity at integer electron numbers: the functional
derivative $v_{xc}[\rho]$ has a finite jump as $N$ passes through
an integer.  This is essential for the correct prediction of
fundamental gaps and is a theorem of exact DFT.

\emph{Step 3: Contradiction.}

A globally smooth $v_{xc}=2\eta\Delta\rho$ cannot exhibit a
derivative discontinuity at any $N$.  Therefore
$E_{xc}[\rho]=\eta\langle\rho,\Delta\rho\rangle$ cannot be the
exact universal functional.
\end{proof}

\begin{remark}[The correct LoNalogy target]
The no-go applies to \emph{universal fixed} kernels.  The
system-dependent bridge covariance functional
$E_{\br,\mathrm{cov}}[X|X_0]$ of Theorem~\ref{thm:functional}
is not fixed: it depends on $X$ through
$\mathcal{C}_{\br}[X_s]$, which varies along the integration
path.  This state-dependence is essential and allows the
functional to be compatible with derivative discontinuities.
\end{remark}


%% ============================================================
\part{Modified P\"oschl--Teller Local Normal Form}
%% ============================================================

\section{Local Surrogate Construction}\label{sec:mPT}

\begin{definition}[Modified P\"oschl--Teller potential]
The modified P\"oschl--Teller (mPT) potential is
\begin{equation}
V_{\mathrm{mPT}}(x;\ell,D,V_\infty)
= V_\infty - D\,\sech^2(x/\ell),
\end{equation}
where $D>0$ is the well depth, $\ell>0$ is the length scale, and
$V_\infty$ is the asymptotic value.
\end{definition}

\begin{theorem}[Well-branch local normal form]\label{thm:well}
Let $V(q)$ be a smooth 1D potential with a non-degenerate local
minimum at $q_*$: $V'(q_*)=0$, $k:=V''(q_*)>0$.
Fix an asymptotic reference energy $V_\infty>V(q_*)$ and define
$D:=V_\infty-V(q_*)>0$ and $\ell:=\sqrt{2D/k}$.

Then there exists a smooth local coordinate change
$\Delta q=q-q_*\mapsto x=x(\Delta q)$ with $x(0)=0$, $x'(0)=1$,
such that
\begin{equation}\label{eq:wellnormal}
\boxed{V(q(x)) = V_\infty - D\,\sech^2(x/\ell) + O(x^5).}
\end{equation}

Explicitly, the coordinate change is
\begin{equation}\label{eq:coordwell}
\Delta q = x
-\frac{g}{6k}\,x^2
+\left(-\frac{k}{6D}+\frac{5g^2}{72k^2}
-\frac{h}{24k}\right)x^3 + O(x^4),
\end{equation}
where $g:=V'''(q_*)$ and $h:=V''''(q_*)$.
\end{theorem}

\begin{proof}
\emph{Step 1: Taylor expand the mPT potential.}

Using $\sech^2 u=1-u^2+\frac{2}{3}u^4+O(u^6)$:
\begin{align}
V_\infty-D\sech^2(x/\ell)
&= V_\infty-D\left(1-\frac{x^2}{\ell^2}
+\frac{2x^4}{3\ell^4}+O(x^6)\right)\\
&= V(q_*)+\frac{D}{\ell^2}x^2
-\frac{2D}{3\ell^4}x^4+O(x^6)\\
&= V(q_*)+\frac{k}{2}x^2-\frac{k^2}{6D}x^4+O(x^6),
\end{align}
where we used $D/\ell^2=k/2$.

\emph{Step 2: Taylor expand $V(q)$ in $\Delta q$.}

\begin{equation}
V(q_*+\Delta q) = V(q_*)+\frac{k}{2}\Delta q^2
+\frac{g}{6}\Delta q^3+\frac{h}{24}\Delta q^4+O(\Delta q^5).
\end{equation}

\emph{Step 3: Ansatz for coordinate change.}

Write $\Delta q=x+\alpha_2 x^2+\alpha_3 x^3+O(x^4)$.
Substituting into the Taylor expansion of $V$:
\begin{align}
V(q(x)) &= V(q_*)+\frac{k}{2}(x+\alpha_2 x^2+\alpha_3 x^3+\cdots)^2\\
&\quad+\frac{g}{6}(x+\alpha_2 x^2+\cdots)^3
+\frac{h}{24}(x+\cdots)^4+O(x^5).
\end{align}

\emph{Step 4: Match coefficients.}

\underline{$x^2$ coefficient:}
$V$-side: $\frac{k}{2}$.  mPT-side: $\frac{k}{2}$.  Matches
automatically.

\underline{$x^3$ coefficient:}
$V$-side: $k\alpha_2+\frac{g}{6}$.
mPT-side: $0$ (the mPT potential is even in $x$, so the cubic
term vanishes).
Setting equal: $\alpha_2=-g/(6k)$.

\underline{$x^4$ coefficient:}
$V$-side: $k\alpha_3+\frac{k}{2}\alpha_2^2
+\frac{g}{2}\alpha_2+\frac{h}{24}$.
mPT-side: $-k^2/(6D)$.
Setting equal:
\begin{equation}
\alpha_3 = -\frac{k}{6D}+\frac{5g^2}{72k^2}-\frac{h}{24k}.
\end{equation}
(Using $\alpha_2=-g/(6k)$:
$\frac{k}{2}\alpha_2^2=\frac{g^2}{72k}$,
$\frac{g}{2}\alpha_2=-\frac{g^2}{12k}$;
then $k\alpha_3=-k^2/(6D)-g^2/(72k)+g^2/(12k)-h/24
=-k^2/(6D)+5g^2/(72k)-h/24$,
so $\alpha_3=-k/(6D)+5g^2/(72k^2)-h/(24k)$.)

This gives~\eqref{eq:coordwell}, and the matching through
$O(x^4)$ gives~\eqref{eq:wellnormal}.
\end{proof}

\begin{theorem}[Barrier-branch local normal form]\label{thm:barrier}
Let $V(q)$ have a non-degenerate local maximum at $q_*$:
$\kappa:=-V''(q_*)>0$.  Fix $V_{\mathrm{base}}<V(q_*)$ and define
$D:=V(q_*)-V_{\mathrm{base}}>0$, $\ell:=\sqrt{2D/\kappa}$.

Then there exists a smooth coordinate change with
\begin{equation}\label{eq:barnormal}
\boxed{V(q(x)) = V_{\mathrm{base}} + D\,\sech^2(x/\ell) + O(x^5),}
\end{equation}
where
\begin{equation}\label{eq:coordbar}
\Delta q = x
+\frac{g}{6\kappa}\,x^2
+\left(-\frac{\kappa}{6D}+\frac{5g^2}{72\kappa^2}
+\frac{h}{24\kappa}\right)x^3 + O(x^4).
\end{equation}
\end{theorem}

\begin{proof}
Identical to the well-branch proof with $k\to-\kappa$,
$V_\infty\to V_{\mathrm{base}}$, and appropriate sign adjustments.
The mPT barrier potential
$V_{\mathrm{base}}+D\sech^2(x/\ell)$ has Taylor expansion
$V(q_*)-\frac{\kappa}{2}x^2+\frac{\kappa^2}{6D}x^4+O(x^6)$,
and the matching proceeds analogously.
\end{proof}


\section{No-Go: Universal PES Law}\label{sec:nogo_PES}

\begin{theorem}[Stationary-point invariant obstruction]\label{thm:nogo_PES}
If $V(q)=V_\infty-D\sech^2((q-q_*)/\ell)$ holds exactly in a
neighbourhood of a local minimum $q_*$ (in the original coordinate
$q$, without coordinate change), then necessarily
\begin{equation}\label{eq:obstruction}
\boxed{V'''(q_*)=0,\qquad
V''''(q_*)=-\frac{4[V''(q_*)]^2}{D}.}
\end{equation}
For the barrier branch,
$V''''(q_*)=+4[V''(q_*)]^2/D$.

Generic chemical potentials violate~\eqref{eq:obstruction}
(e.g., the Morse potential has $V'''(q_*)\neq 0$).  Therefore:
\begin{equation}
\boxed{\text{Modified P\"oschl--Teller is not a universal PES law.}}
\end{equation}
\end{theorem}

\begin{proof}
Expanding $\sech^2((q-q_*)/\ell)$ directly in $q$:
\begin{align}
V_\infty-D\sech^2\frac{q-q_*}{\ell}
&= V(q_*)+\frac{D}{\ell^2}(q-q_*)^2
-\frac{2D}{3\ell^4}(q-q_*)^4+O((q-q_*)^6).
\end{align}
(The $\sech^2$ function is even, so all odd-order terms vanish.)

If this equals $V(q)=V(q_*)+\frac{k}{2}(q-q_*)^2
+\frac{g}{6}(q-q_*)^3+\frac{h}{24}(q-q_*)^4+\cdots$,
then:

\underline{$O((q-q_*)^3)$:} $g/6=0$, so $V'''(q_*)=g=0$.

\underline{$O((q-q_*)^4)$:} $h/24=-2D/(3\ell^4)$.
Using $\ell^2=2D/k$: $-2D/(3\ell^4)=-2D/(3\cdot4D^2/k^2)
=-k^2/(6D)$.
So $h=-4k^2/D$, i.e., $V''''(q_*)=-4[V''(q_*)]^2/D$.

For the Morse potential: $V_M'''(q_*)=-Da^3\neq 0$ generically.
Therefore the Morse potential violates $V'''(q_*)=0$.
\end{proof}

\begin{corollary}[Error order]\label{cor:error}
For a generic potential with $V'''(q_*)\neq 0$, the error of the
mPT surrogate \emph{without} coordinate change is $O((q-q_*)^3)$.
With the canonical coordinate change of
Theorem~\ref{thm:well}/\ref{thm:barrier}, the error is $O(x^5)$,
i.e., fourth-order agreement.
\end{corollary}


%% ============================================================
\part{Complete Classification}
%% ============================================================

\section{Summary}\label{sec:summary}

\begin{center}
\begin{tabular}{p{5.5cm}lll}
\toprule
Result & Ref & Status\\
\midrule
\multicolumn{3}{l}{\textbf{(A) Active-space downfolding}}\\
\quad Canonical bridge selection & Thm~\ref{thm:decoupling} & Exact\\
\quad Bridge rank bound & Prop~\ref{prop:rankbound} & Exact\\
\quad Three-sector Feshbach--Schur & Thm~\ref{thm:downfolding} & Exact\\
\quad Fixed-point iteration & Thm~\ref{thm:fixedpoint} & Exact\\
\quad Static generalised eigenproblem & Thm~\ref{thm:static} & Exact\\
\quad Energy derivative & Prop~\ref{prop:Rstar} & Exact\\
\midrule
\multicolumn{3}{l}{\textbf{(B) Bridge covariance functional}}\\
\quad Log-marginal Hessian & Thm~\ref{thm:hessian} & Exact\\
\quad Path-integrated functional & Thm~\ref{thm:functional} & Exact\\
\quad Gaussian bridge kernel & Thm~\ref{thm:gaussian} & Exact\\
\quad No-go: fixed $E_{xc}$ & Thm~\ref{thm:nogo_Exc} & No-go\\
\midrule
\multicolumn{3}{l}{\textbf{(C) mPT local normal form}}\\
\quad Well branch (4th order) & Thm~\ref{thm:well} & Exact\\
\quad Barrier branch (4th order) & Thm~\ref{thm:barrier} & Exact\\
\quad No-go: universal PES law & Thm~\ref{thm:nogo_PES} & No-go\\
\quad Error order & Cor~\ref{cor:error} & Exact\\
\bottomrule
\end{tabular}
\end{center}

\section{Final Verdict}\label{sec:verdict}

The LoNalogy toolkit for quantum chemistry consists of:
\begin{equation}
\boxed{\text{Canonical bridge}\;+\;\text{Exact Feshbach--Schur}
\;+\;\text{Bridge covariance}\;+\;\text{Path-specific mPT.}}
\end{equation}

What is closed:
\begin{enumerate}[label=(\Alph*)]
\item Active-space downfolding is the exact algebraic parent of
CASPT2/NEVPT2/DUCC.
\item The exact correlation correction is a system-dependent bridge
covariance functional, not a universal fixed kernel.
\item The mPT potential is a canonical local normal form at
stationary points, with fourth-order agreement under coordinate
change.
\end{enumerate}

What is provably impossible (no-go):
\begin{enumerate}[label=(\alph*)]
\item $\sech^2$ as a universal chemical bond law (cubic obstruction).
\item Fixed quadratic $E_{xc}$ as the exact universal functional
(derivative discontinuity).
\end{enumerate}

What remains for implementation:
\begin{quote}
Numerical benchmarks on LiH, H$_6$, H$_2$O bond dissociation,
comparing LoNalogy exact downfolding against CASPT2, NEVPT2, and
DUCC.
\end{quote}


\section*{Conclusion}
\addcontentsline{toc}{section}{Conclusion}

LoNalogy does not change quantum chemistry by replacing string
theory.  It changes quantum chemistry by providing exact algebraic
tools for multi-electron correlation: a canonical bridge selection
that eliminates the arbitrariness of active-space partitioning, an
exact Feshbach--Schur downfolding that subsumes all perturbative
and coupled-cluster embedding methods as approximations, a
system-dependent bridge covariance functional that identifies the
correct target for correlation corrections (not a fixed kernel),
and a path-specific $\sech^2$ surrogate that provides analytically
solvable local models of reaction barriers.

In one sentence:
\begin{quote}
\emph{LoNalogy's quantum chemistry core is:
canonical bridge $+$ exact Schur downfolding $+$
bridge covariance correction $+$ path-specific mPT normal form.}
\end{quote}


\section*{Acknowledgements}
\addcontentsline{toc}{section}{Acknowledgements}

The author thanks Claude (Anthropic) for collaborative development
of the bridge covariance formalism and verification of all proofs.


\begin{thebibliography}{99}

\bibitem{LoNalogyv87}
M.~Nakamura,
\textit{LoNalogy v87.1},
Zenodo (2026),
\href{https://doi.org/10.5281/zenodo.19115338}{doi:10.5281/zenodo.19115338}.

\bibitem{YMmassGap}
M.~Nakamura,
\textit{Yang--Mills Existence and Mass Gap},
Zenodo (2026),
\href{https://doi.org/10.5281/zenodo.19183838}{doi:10.5281/zenodo.19183838}.

\bibitem{Feshbach}
H.~Feshbach,
\textit{Unified theory of nuclear reactions},
Ann.\ Phys.\ \textbf{5}, 357 (1958).

\bibitem{CASPT2}
K.~Andersson, P.-\AA.~Malmqvist, and B.~O.~Roos,
\textit{Second-order perturbation theory with a complete active
space self-consistent field reference function},
J.\ Chem.\ Phys.\ \textbf{96}, 1218 (1992).

\bibitem{DUCC}
N.~P.~Bauman \textit{et al.},
\textit{Downfolding of many-body Hamiltonians using
active-space models},
J.\ Chem.\ Phys.\ \textbf{151}, 014107 (2019).

\bibitem{PoschlTeller}
G.~P\"oschl and E.~Teller,
\textit{Bemerkungen zur Quantenmechanik des anharmonischen
Oszillators},
Z.\ Phys.\ \textbf{83}, 143 (1933).

\bibitem{KohnSham}
W.~Kohn and L.~J.~Sham,
\textit{Self-Consistent Equations Including Exchange and
Correlation Effects},
Phys.\ Rev.\ \textbf{140}, A1133 (1965).

\end{thebibliography}

\end{document}
